\documentclass[12pt,a4paper]{article}

% =====================================================================
%  PACOTES ESSENCIAIS
% =====================================================================
\usepackage[brazil]{babel}
\usepackage[utf8]{inputenc}
\usepackage[T1]{fontenc}
\usepackage{graphicx}
\usepackage{cite}
\usepackage{amsmath}
\usepackage{amssymb}
\usepackage{indentfirst}
\usepackage{geometry}
\usepackage{booktabs}
\usepackage{multirow}
\usepackage{array}
\usepackage{subcaption}
\usepackage{float}
\usepackage{listings}
\usepackage{xcolor}
\usepackage{tikz}
\usepackage{tcolorbox}
\usepackage{fancyhdr}
\usepackage{titlesec}
\usepackage{enumitem}
\usepackage{pgfplots}
\usepackage{longtable}
\usepackage{tabularx}
\usepackage{microtype}
\usepackage{url}
\usepackage{hyperref}

% =====================================================================
%  CONFIGURACAO PGFPLOTS / TIKZ
% =====================================================================
\pgfplotsset{compat=1.18}
\usetikzlibrary{shapes.geometric, arrows, positioning, fit}

% =====================================================================
%  CORES
% =====================================================================
\definecolor{primaryblue}{RGB}{41, 128, 185}
\definecolor{accentblue}{RGB}{52, 152, 219}
\definecolor{successgreen}{RGB}{39, 174, 96}
\definecolor{warningorange}{RGB}{243, 156, 18}
\definecolor{dangerred}{RGB}{231, 76, 60}
\definecolor{lightgray}{RGB}{236, 240, 241}
\definecolor{mediumgray}{RGB}{189, 195, 199}
\definecolor{darkgray}{RGB}{44, 62, 80}
\definecolor{purpleaccent}{RGB}{142, 68, 173}
\definecolor{deepblue}{RGB}{25, 42, 86}
\definecolor{codegreen}{rgb}{0,0.6,0}
\definecolor{codegray}{rgb}{0.5,0.5,0.5}
\definecolor{codepurple}{rgb}{0.58,0,0.82}
\definecolor{backcolour}{rgb}{0.97,0.97,0.97}

% =====================================================================
%  LAYOUT
% =====================================================================
\geometry{a4paper,left=3cm,right=2cm,top=3cm,bottom=2cm}

% =====================================================================
%  CABECALHO E RODAPE
% =====================================================================
\pagestyle{fancy}
\fancyhf{}
\fancyhead[L]{\small\textcolor{darkgray}{Detec\c{c}\~{a}o de \textit{deepfakes} de voz}}
\fancyhead[R]{\small\textcolor{darkgray}{T.~G.~Braga}}
\fancyfoot[C]{\textcolor{darkgray}{\thepage}}
\renewcommand{\headrulewidth}{0.5pt}
\renewcommand{\footrulewidth}{0pt}

% =====================================================================
%  TITULOS
% =====================================================================
\titleformat{\section}
  {\Large\bfseries\color{deepblue}}{\thesection}{1em}{}
\titleformat{\subsection}
  {\large\bfseries\color{darkgray}}{\thesubsection}{1em}{}
\titleformat{\subsubsection}
  {\normalsize\bfseries\color{darkgray}}{\thesubsubsection}{1em}{}

% =====================================================================
%  CAIXAS INFORMATIVAS
% =====================================================================
\tcbuselibrary{skins,breakable}

\newtcolorbox{infobox}[1]{
    colback=primaryblue!5, colframe=primaryblue,
    coltitle=white, fonttitle=\bfseries,
    title=#1, rounded corners, breakable}

\newtcolorbox{warningbox}[1]{
    colback=warningorange!10, colframe=warningorange,
    coltitle=white, fonttitle=\bfseries,
    title=#1, rounded corners}

\newtcolorbox{successbox}[1]{
    colback=successgreen!10, colframe=successgreen,
    coltitle=white, fonttitle=\bfseries,
    title=#1, rounded corners}

% =====================================================================
%  LISTAS
% =====================================================================
\setlist[itemize]{label=\textcolor{primaryblue}{\textbullet}, leftmargin=1.5em}
\setlist[enumerate]{label=\textcolor{primaryblue}{\arabic*.}, leftmargin=1.5em}

% =====================================================================
%  CODIGO
% =====================================================================
\lstdefinestyle{modernstyle}{
    backgroundcolor=\color{backcolour},
    commentstyle=\color{codegreen},
    keywordstyle=\color{purpleaccent}\bfseries,
    numberstyle=\tiny\color{codegray},
    stringstyle=\color{codepurple},
    basicstyle=\ttfamily\small,
    breaklines=true,
    captionpos=b,
    keepspaces=true,
    numbers=left,
    numbersep=8pt,
    showstringspaces=false,
    tabsize=2,
    frame=single,
    framerule=0.5pt,
    rulecolor=\color{mediumgray}
}
\lstset{style=modernstyle}

% =====================================================================
%  HYPERLINKS
% =====================================================================
\hypersetup{
    colorlinks=true,
    linkcolor=primaryblue,
    filecolor=purpleaccent,
    urlcolor=primaryblue,
    citecolor=primaryblue,
    bookmarksopen=true,
    pdfstartview=FitH,
    pdftitle={Pipeline Modular para Detecção de Áudio Sintético},
    pdfauthor={Thierry Gonçalves Braga},
    pdfsubject={Trabalho de Conclusão de Curso},
    pdfkeywords={deepfake, áudio, detecção, machine learning}
}

% =====================================================================
%  ESTILOS TIKZ
% =====================================================================
\tikzset{
  block/.style={rectangle, rounded corners=5pt, minimum width=3.2cm,
    minimum height=1cm, text centered, draw=primaryblue, line width=1.2pt,
    fill=primaryblue!10, font=\small\sffamily},
  decision/.style={diamond, minimum width=2.5cm, minimum height=1cm,
    text centered, draw=warningorange, line width=1.2pt,
    fill=warningorange!15, font=\small\sffamily},
  startend/.style={rectangle, rounded corners=8pt, minimum width=3cm,
    minimum height=0.9cm, text centered, draw=successgreen, line width=1.5pt,
    fill=successgreen!20, font=\small\sffamily\bfseries},
  modelnode/.style={rectangle, rounded corners=6pt, minimum width=2.8cm,
    minimum height=0.9cm, text centered, draw=purpleaccent, line width=1.5pt,
    fill=purpleaccent!15, font=\small\sffamily\bfseries},
  resultreal/.style={rectangle, rounded corners=6pt, minimum width=2.5cm,
    minimum height=0.9cm, text centered, draw=successgreen, line width=1.5pt,
    fill=successgreen!25, font=\small\sffamily\bfseries},
  resultfake/.style={rectangle, rounded corners=6pt, minimum width=2.5cm,
    minimum height=0.9cm, text centered, draw=dangerred, line width=1.5pt,
    fill=dangerred!25, font=\small\sffamily\bfseries},
  arrow/.style={thick,->,>=stealth,color=darkgray,line width=1.3pt}
}

% =====================================================================
%  DOCUMENTO
% =====================================================================
\begin{document}
\pagestyle{empty}

% ------------------------------------------------------------------
%  PAGINA DE TITULO (NBR 14724:2011)
% ------------------------------------------------------------------
\begin{center}
{\large UNIVERSIDADE FEDERAL DE S\~{A}O JO\~{A}O DEL-REI}\\
{\normalsize Departamento de Ci\^{e}ncia da Computa\c{c}\~{a}o}\\[2.0cm]

{\Large\bfseries
PIPELINE MODULAR PARA DETEC\c{C}\~{A}O DE \'{A}UDIO SINT\'{E}TICO\\[0.3cm]
USANDO ARQUITETURAS NEURAIS H\'{I}BRIDAS E\\[0.3cm]
AN\'{A}LISE EMP\'{I}RICA DE CARACTER\'{I}STICAS}\\[2.0cm]

{\large Thierry Gon\c{c}alves Braga}\\[0.5cm]

{\normalsize
Trabalho de Conclus\~{a}o de Curso apresentado \`{a} Universidade\\
Federal de S\~{a}o Jo\~{a}o del-Rei}\\[0.5cm]

{\normalsize Orientadora: Prof.\textsuperscript{a}\ Ana Claudia}\\[3.0cm]

{\normalsize S\~{a}o Jo\~{a}o del-Rei --- MG\\
2026}
\end{center}

\newpage

% ------------------------------------------------------------------
%  FOLHA DE APROVAÇÃO
% ------------------------------------------------------------------
\begin{center}
\textbf{FOLHA DE APROVA\c{C}\~{A}O}
\end{center}

\vspace{1.5cm}
\noindent
Trabalho de Conclus\~{a}o de Curso defendido e aprovado em \underline{\hspace{4cm}},
como requisito parcial para obten\c{c}\~{a}o do t\'{i}tulo de Bacharel em
Ci\^{e}ncia da Computa\c{c}\~{a}o.

\vspace{2cm}
\noindent\textbf{Banca Examinadora:}

\vspace{1cm}
\noindent\rule{10cm}{0.4pt}\\
Prof.\textsuperscript{a}\ Ana Claudia (Orientadora)\\
Universidade Federal de S\~{a}o Jo\~{a}o del-Rei

\vspace{0.8cm}
\noindent\rule{10cm}{0.4pt}\\
Prof.\ [Nome do Membro 1]\\
[Institui\c{c}\~{a}o]

\vspace{0.8cm}
\noindent\rule{10cm}{0.4pt}\\
Prof.\ [Nome do Membro 2]\\
[Institui\c{c}\~{a}o]

\newpage

% ------------------------------------------------------------------
%  DEDICATÓRIA
% ------------------------------------------------------------------
\begin{center}\textit{Dedicat\'{o}ria}\end{center}
\vfill
\begin{flushright}
\textit{\`{A} minha fam\'{i}lia, pelo apoio incondicional.}
\end{flushright}
\vfill
\newpage

% ------------------------------------------------------------------
%  AGRADECIMENTOS
% ------------------------------------------------------------------
\begin{center}\textbf{AGRADECIMENTOS}\end{center}
\noindent
Agrade\c{c}o \`{a} minha orientadora, Prof.\textsuperscript{a}\ Ana Claudia,
pela orienta\c{c}\~{a}o t\'{e}cnica e cient\'{i}fica.
Aos colegas do laborat\'{o}rio, pelas discuss\~{o}es enriquecedoras.
\`{A} UFSJ, pela infraestrutura.
\newpage

% ------------------------------------------------------------------
%  RESUMO (NBR 6028:2021)
% ------------------------------------------------------------------
\begin{center}\textbf{RESUMO}\end{center}

\noindent
Este trabalho apresenta o desenvolvimento de um \textit{pipeline} modular e de
c\'{o}digo aberto para detec\c{c}\~{a}o de \'{a}udio sint\'{e}tico, abordando
o desafio crescente da prolifera\c{c}\~{a}o de \textit{deepfakes} de voz.
A metodologia integra pr\'{e}-processamento robusto com detec\c{c}\~{a}o de
atividade de voz (VAD) e normaliza\c{c}\~{a}o autom\'{a}tica de ganho (AGC),
extra\c{c}\~{a}o emp\'{i}rica de caracter\'{i}sticas ac\'{u}sticas ranqueadas
por desempenho discriminativo e classifica\c{c}\~{a}o por meio de cinco
fam\'{i}lias de arquiteturas neurais implementadas em TensorFlow~2.21/Keras~3.
Um \textit{survey} sistem\'{a}tico das gera\c{c}\~{o}es de m\'{e}todos de
detec\c{c}\~{a}o contextualiza o posicionamento do sistema.
Os experimentos foram conduzidos em CPU (sem GPU) sobre um \textit{dataset}
balanceado de 600 amostras de fala portuguesas (BRSpeech-DF), com divis\~{a}o
70/15/15.
Os resultados revelam uma importante dicotomia pr\'{a}tica: arquiteturas
baseadas em espectrograma convergiram com sucesso, enquanto modelos de
an\'{a}lise direta de forma de onda (raw-audio) n\~{a}o convergiram nas
condi\c{c}\~{o}es experimentais, resultado consistente com os requisitos
computacionais documentados na literatura.
A MultiscaleCNN (Res2Net-50) alcan\c{c}ou 97,78\% de acur\'{a}cia e EER de
5,56\%, com AUC-ROC de 0,997.
O Ensemble Adaptativo obteve 91,11\% de acur\'{a}cia, EER de 13,33\% e
lat\^{e}ncia de 48,1\,ms.
Os testes de robustez sob ru\'{i}do AWGN revelaram degrada\c{c}\~{a}o
substancial abaixo de SNR~30\,dB, evidenciando a necessidade de
\textit{data augmentation} com ru\'{i}do para uso em campo.
O sistema \'{e} integralmente constru\'{i}do sobre tecnologias livres
(TensorFlow, Keras, Librosa), garantindo reprodutibilidade e extensibilidade.

\vspace{0.4cm}
\noindent\textbf{Palavras-chave:} \textit{deepfake} de \'{a}udio.
\textit{pipeline} modular. Arquiteturas neurais h\'{i}bridas.
MultiscaleCNN. Ensemble adaptativo. Robustez a ru\'{i}do. C\'{o}digo aberto.

\vspace{1cm}
\begin{center}\textbf{ABSTRACT}\end{center}

\noindent
This work presents the development of a modular open-source pipeline for
synthetic audio detection, addressing the growing challenge of voice deepfake
proliferation.
The methodology integrates robust preprocessing with Voice Activity Detection
(VAD) and Automatic Gain Control (AGC), empirical acoustic feature extraction
ranked by discriminative performance, and classification through five families
of neural architectures implemented in TensorFlow~2.21/Keras~3.
Experiments were conducted on CPU (no GPU) over a balanced 600-sample
Portuguese speech dataset (BRSpeech-DF), with 70/15/15 stratified split.
Results reveal an important practical dichotomy: spectrogram-based architectures
converged successfully, while direct raw-waveform models did not converge
under the experimental conditions, a result consistent with computational
requirements documented in the literature.
The MultiscaleCNN (Res2Net-50) achieved 97.78\% accuracy and 5.56\% EER,
with AUC-ROC of 0.997.
The Adaptive Ensemble achieved 91.11\% accuracy, 13.33\% EER, and 48.1\,ms
latency.
Robustness tests under AWGN noise revealed substantial degradation below
SNR~30\,dB, highlighting the need for noise-based data augmentation for
deployment.
The system is fully built on open-source technologies (TensorFlow, Keras,
Librosa), ensuring reproducibility and extensibility.

\vspace{0.4cm}
\noindent\textbf{Keywords:} Audio deepfake. Modular pipeline.
Hybrid neural architectures. MultiscaleCNN. Adaptive ensemble.
Noise robustness. Open-source.

\newpage
\pagestyle{fancy}
\tableofcontents
\newpage

% =====================================================================
%  1. INTRODUCAO
% =====================================================================
\section{Introdu\c{c}\~{a}o}

A crescente democratiza\c{c}\~{a}o de tecnologias de s\'{i}ntese de \'{a}udio
baseadas em intelig\^{e}ncia artificial tem facilitado a cria\c{c}\~{a}o de
\textit{deepfakes} de \'{a}udio com qualidade cada vez mais pr\'{o}xima \`{a}
voz humana natural \cite{muller2022warning, heidari2024deepfake}.
T\'{e}cnicas modernas como RVC (\textit{Retrieval-based Voice Conversion})
\cite{bird2023realtime}, Tacotron~2 \cite{shen2018natural} e WaveNet
\cite{oord2016wavenet} permitem a gera\c{c}\~{a}o de conte\'{u}do sint\'{e}tico
realista com apenas alguns minutos de \'{a}udio original.

Esta evolu\c{c}\~{a}o tecnol\'{o}gica, embora ben\'{e}fica para aplica\c{c}\~{o}es
leg\'{i}timas como assistentes virtuais e dublagem automatizada, tem sido
explorada maliciosamente em fraudes financeiras, manipula\c{c}\~{a}o pol\'{i}tica
e chantagem \cite{pindrop2025, dfrlab2024}.
No contexto brasileiro, foram documentados 79 casos de conte\'{u}do sint\'{e}tico
relacionado a candidatos durante as elei\c{c}\~{o}es municipais de 2024
\cite{farrugia2024}, evidenciando a necessidade urgente de sistemas de
detec\c{c}\~{a}o robustos.

Estudos perceptuais indicam que humanos identificam corretamente
\textit{deepfakes} de \'{a}udio em apenas 73\% dos casos, taxa que decresce
para 51\% com \'{a}udios de alta qualidade \cite{muller2022warning}.
Isso evidencia a necessidade de sistemas automatizados robustos e explica o
surgimento dos \textit{benchmarks} padronizados ASVspoof como eixo central de
avalia\c{c}\~{a}o \cite{todisco2019asvspoof, yamagishi2021asvspoof,
wang2024asvspoof5}.

\subsection{Objetivos}

\subsubsection{Objetivo Geral}
Desenvolver um \textit{pipeline} modular, de c\'{o}digo aberto e extens\'{i}vel
para detec\c{c}\~{a}o de \'{a}udio sint\'{e}tico que integre m\'{u}ltiplas
arquiteturas neurais especializadas e permita an\'{a}lise emp\'{i}rica
sistem\'{a}tica de caracter\'{i}sticas ac\'{u}sticas.

\subsubsection{Objetivos Espec\'{i}ficos}
\begin{itemize}
    \item Realizar um \textit{survey} das gera\c{c}\~{o}es de m\'{e}todos de
    detec\c{c}\~{a}o de \'{a}udio sint\'{e}tico.
    \item Implementar pr\'{e}-processamento com VAD e AGC.
    \item Ranquear caracter\'{i}sticas ac\'{u}sticas por capacidade
    discriminativa, ancorado na literatura consolidada.
    \item Implementar e comparar cinco arquiteturas neurais especializadas
    sobre um \textit{dataset} brasileiro balanceado.
    \item Avaliar robustez sob ru\'{i}do AWGN com m\'{u}ltiplos n\'{i}veis
    de SNR.
    \item Propor um \textit{pipeline} reprodut\'{i}vel com tecnologias de
    c\'{o}digo aberto.
\end{itemize}

% =====================================================================
%  2. FUNDAMENTACAO TEORICA
% =====================================================================
\section{Fundamenta\c{c}\~{a}o Te\'{o}rica}
\label{sec:fundamentacao}

\subsection{Tecnologias de S\'{i}ntese de Voz}

\subsubsection{Clonagem de Voz}

A clonagem de voz (\textit{voice cloning}) consiste em criar um modelo capaz de
replicar as caracter\'{i}sticas ac\'{u}sticas e pros\'{o}dicas de um falante
espec\'{i}fico a partir de amostras limitadas \cite{muller2022warning}.
Modelos como VALL-E \cite{wang2023vall} utilizam abordagem \textit{few-shot}
baseada em codificadores neurais de \'{a}udio, atingindo qualidade quase
indistingu\'{i}vel com apenas 3 segundos de \'{a}udio.
O processo pode ser formalizado como:

\begin{equation}
\mathbf{y}_{clone} = f_{\theta}(\mathbf{x}_{text},\; \mathbf{z}_{speaker})
\end{equation}

onde $\mathbf{y}_{clone}$ \'{e} o \'{a}udio sint\'{e}tico gerado,
$\mathbf{x}_{text}$ \'{e} o texto de entrada, $\mathbf{z}_{speaker}$ \'{e} o
vetor de caracter\'{i}sticas do falante e $f_{\theta}$ \'{e} o modelo neural
treinado.

\begin{infobox}{Clonagem de Voz em Contextos Maliciosos}
A clonagem de voz foi documentada em fraudes financeiras, com casos reportados
em 2024 onde \'{a}udios sint\'{e}ticos imitando executivos induziram
funcion\'{a}rios a realizar transfer\^{e}ncias de alto valor. A acessibilidade
de ferramentas de c\'{o}digo aberto amplifica significativamente esse risco
\cite{pindrop2025}.
\end{infobox}

\subsubsection{Convers\~{a}o de Voz}

A convers\~{a}o de voz (\textit{voice conversion}, VC) transforma a voz de um
falante para soar como a de outro, preservando o conte\'{u}do lingu\'{i}stico
original \cite{bird2023realtime}.
A t\'{e}cnica RVC emprega bancos de dados de vozes e extratores de
caracter\'{i}sticas para mapear o falante de origem ao alvo:

\begin{equation}
\mathbf{y}_{conv} = g_{\phi}(\mathbf{x}_{source},\; \mathbf{z}_{target})
\end{equation}

Modelos como StarGAN-VC \cite{kameoka2018stargan} utilizam redes adversariais
para aprender mapeamentos muitos-para-muitos entre vozes com lat\^{e}ncia
reduzida.
A VC \'{e} desafiadora para detec\c{c}\~{a}o, pois preserva o conte\'{u}do
lingu\'{i}stico, mascarando artefatos e levando detectores baseados em
caracter\'{i}sticas espectrais a EERs de at\'{e} 15\%
\cite{wang2024asvspoof5}.

\begin{infobox}{Convers\~{a}o de Voz nas Elei\c{c}\~{o}es Brasileiras de 2024}
Durante as elei\c{c}\~{o}es municipais brasileiras de 2024, a convers\~{a}o de
voz foi identificada em casos de manipula\c{c}\~{a}o de discursos de
candidatos, onde trechos de \'{a}udio foram alterados para simular falas
comprometedoras \cite{farrugia2024}.
\end{infobox}

\subsubsection{S\'{i}ntese de Fala (\textit{Text-to-Speech})}

O Tacotron~2 \cite{shen2018natural} estabeleceu arquiteturas
\textit{sequence-to-sequence} para mapear texto em espectrogramas mel,
convertidos em \'{a}udio por um vocoder como WaveNet \cite{oord2016wavenet}:

\begin{equation}
\mathbf{s}_{mel} = h_{\psi}(\mathbf{t}_{text}), \qquad
\mathbf{y}_{audio} = v_{\omega}(\mathbf{s}_{mel})
\end{equation}

A WaveNet modela diretamente a forma de onda bruta, alcan\c{c}ando qualidade
pr\'{o}xima \`{a} humana \cite{oord2016wavenet}. Sistemas TTS avan\c{c}ados
dificultam a detec\c{c}\~{a}o por m\'{e}todos tradicionais baseados em MFCCs
devido \`{a} alta naturalidade dos \'{a}udios gerados.

\subsubsection{Manipula\c{c}\~{a}o em Tempo Real}

A manipula\c{c}\~{a}o em tempo real permite alterar \'{a}udios durante
conversas ao vivo com lat\^{e}ncia inferior a 200\,ms, tornando-a
impercept\'{i}vel para humanos \cite{pindrop2025}.
Sistemas como RVC~v2 otimizam a convers\~{a}o para baixa lat\^{e}ncia
utilizando redes convolucionais 1D leves \cite{bird2023realtime}:

\begin{equation}
\mathbf{y}_{rt} = r_{\xi}(\mathbf{x}_{input},\; \mathbf{z}_{target},\; t)
\end{equation}

\begin{warningbox}{Riscos do \textit{vishing} em Tempo Real}
Casos de \textit{vishing} em 2025 usaram manipula\c{c}\~{a}o em tempo real
para imitar vozes de familiares, enganando v\'{i}timas em transfer\^{e}ncias
financeiras de emerg\^{e}ncia \cite{pindrop2025}. A lat\^{e}ncia m\'{e}dia
desses sistemas \'{e} inferior a 200\,ms.
\end{warningbox}

% ------------------------------------------------------------------
\subsection{Gera\c{c}\~{o}es de M\'{e}todos de Detec\c{c}\~{a}o}
% ------------------------------------------------------------------

\subsubsection{Primeira Gera\c{c}\~{a}o --- Caracter\'{i}sticas Manuais e
Estat\'{i}stica Cl\'{a}ssica}

At\'{e} o ASVspoof 2015 \cite{todisco2019asvspoof}, o padr\~{a}o envolvia
extra\c{c}\~{a}o manual de MFCC, LFCC e CQCC seguida por classificadores
GMM-UBM e SVM.
Sahidullah et al.\ \cite{sahidullah2015comparison} demonstraram que o LFCC
supera o MFCC como \textit{front-end} para detec\c{c}\~{a}o de
\textit{spoofing}, enquanto Todisco et al.\ \cite{todisco2016new} introduziram
o CQCC como caracter\'{i}stica de refer\^{e}ncia para a \'{a}rea.
Embora computacionalmente leves, essas arquiteturas falham gravemente com
m\'{e}todos de s\'{i}ntese neural modernos por n\~{a}o generalizarem entre
\textit{datasets} (EER $>$ 20\% em cen\'{a}rios \textit{cross-dataset}).

\subsubsection{Segunda Gera\c{c}\~{a}o --- CNNs e RNNs em Espectrogramas}

A partir do ASVspoof 2019 \cite{todisco2019asvspoof}, CNNs (como ResNet e
Res2Net \cite{gao2021res2net}) e LSTMs processando espectrogramas mel tornaram-se
dominantes.
Sistemas como o LCNN (\textit{Light CNN}) reduziram drasticamente os EERs ao
explorar artefatos visuais de alta frequ\^{e}ncia deixados por
\textit{vocoders}. Abordagens h\'{i}bridas CNN-LSTM reportam acur\'{a}cias
superiores a 90\% \cite{kurnaz2024hybrid}.

\subsubsection{Terceira Gera\c{c}\~{a}o --- An\'{a}lise Direta de Forma de
Onda e Grafos}

A introdu\c{c}\~{a}o de RawNet2 \cite{tak2021rawnet2} e AASIST
\cite{jung2022aasist} marcou o processamento \textit{end-to-end}.
O RawNet2 utiliza \textit{Sinc-Convolutions} para processar diretamente a
forma de onda PCM 1D.
O AASIST implementa Redes de Aten\c{c}\~{a}o em Grafos (GATs) para modelar
depend\^{e}ncias temporais e espectrais simultaneamente, alcan\c{c}ando EER
inferior a 1\% em \textit{datasets} controlados.

\begin{warningbox}{Requisitos Computacionais da 3\textordfeminine\ Gera\c{c}\~{a}o}
Modelos \textit{raw-audio} como RawNet2 e AASIST exigem GPU dedicada e
\textit{datasets} com mais de 5.000 amostras para convergir. Em CPU, mesmo
com 420 amostras de treino, esses modelos n\~{a}o ultrapassam acur\'{a}cia de
acaso (50\%), conforme observado empiricamente neste trabalho e reportado na
literatura \cite{tak2021rawnet2, jung2022aasist}.
\end{warningbox}

\subsubsection{Quarta Gera\c{c}\~{a}o --- Modelos \textit{Foundation} e SSL}

Com o ASVspoof 5 \cite{wang2024asvspoof5} e avalia\c{c}\~{o}es recentes de
modelos autossupervisionados, o padr\~{a}o ouro consolidou-se em torno de
WavLM \cite{chen2022wavlm} e HuBERT \cite{hsu2021hubert}.
Pr\'{e}-treinados em dezenas de milhares de horas de \'{a}udio n\~{a}o
rotulado, esses modelos s\~{a}o ajustados (\textit{fine-tuned}) para
\textit{anti-spoofing}, oferecendo resist\^{e}ncia superior a ataques
adversariais e dados \textit{crowdsourced}.
O Audio Spectrogram Transformer (AST) \cite{gong2021ast} emerge como
alternativa promissora, processando espectrogramas como sequ\^{e}ncias de
\textit{patches} sem convolu\c{c}\~{o}es.

A Tabela~\ref{tab:survey_geracoes} consolida as quatro gera\c{c}\~{o}es de
m\'{e}todos.

\begin{table}[H]
\centering
\caption{\textit{Survey} das Gera\c{c}\~{o}es de M\'{e}todos de Detec\c{c}\~{a}o
de \'{A}udio Sint\'{e}tico}
\label{tab:survey_geracoes}
\small
\begin{tabular}{p{2.0cm}p{3.0cm}p{3.0cm}p{3.0cm}p{2.5cm}}
\toprule
\textbf{Gera\c{c}\~{a}o} & \textbf{Representa\c{c}\~{a}o} &
\textbf{Classificador} & \textbf{Pontos Fortes} & \textbf{Limita\c{c}\~{o}es} \\
\midrule
1\textordfeminine\ (at\'{e} 2015)
  & MFCC, LFCC, CQCC & GMM-UBM, SVM
  & Leveza computacional & Baixa generaliza\c{c}\~{a}o \\
\addlinespace
2\textordfeminine\ (2019)
  & Espectrograma mel & CNN, LSTM, LCNN
  & Robustez a artefatos de vocoder
  & Sens\'{i}vel a geradores n\~{a}o vistos \\
\addlinespace
3\textordfeminine\ (2021)
  & Forma de onda bruta, grafos & RawNet2, AASIST
  & EER $<$ 1\% controlado & GPU + $>$5.000 amostras \\
\addlinespace
4\textordfeminine\ (2024+)
  & Representa\c{c}\~{a}o SSL & WavLM, HuBERT + MLP
  & Generaliza\c{c}\~{a}o \textit{cross-dataset}
  & Depend\^{e}ncia de GPUs potentes \\
\bottomrule
\end{tabular}
\end{table}

% =====================================================================
%  3. METODOLOGIA
% =====================================================================
\section{Metodologia}
\label{sec:metodologia}

\subsection{Vis\~{a}o Geral do Pipeline}

O sistema proposto opera em duas fases: \textbf{treinamento} e
\textbf{predi\c{c}\~{a}o}.
Durante o treinamento, o modelo aprende a distinguir caracter\'{i}sticas de
voz real e sint\'{e}tica.
Na predi\c{c}\~{a}o, processa novos \'{a}udios e fornece classifica\c{c}\~{a}o
bin\'{a}ria (Real/Fake) com \textit{score} de confian\c{c}a.

\begin{figure}[H]
\centering
\begin{tikzpicture}[node distance=1.5cm, scale=0.85, transform shape,
                    every node/.append style={font=\footnotesize}]

% === FASE PREDICAO (esquerda) ===
\draw[primaryblue, line width=1.5pt, rounded corners=6pt, fill=primaryblue!4]
  (-6.8,6.8) rectangle (-0.5,-8.2);
\node[font=\footnotesize\bfseries, color=primaryblue] at (-3.65,6.3)
  {PREDI\c{C}\~{A}O};

\node (audio)   [startend]  at (-3.65, 5.5) {\'{A}udio de Entrada};
\node (vad)     [block, below=of audio]  {VAD + AGC};
\node (feat)    [block, below=of vad]    {Extra\c{c}\~{a}o de Features};
\node (norm)    [block, below=of feat]   {Normaliza\c{c}\~{a}o};
\node (inf)     [modelnode, below=of norm] {Infer\^{e}ncia};
\node (score)   [block, below=of inf]    {Score 0--1};
\node (dec)     [decision, below=of score, yshift=-0.3cm] {$> 0{,}5$?};
\node (real)    [resultreal, below left=1.2cm and 1.0cm of dec]  {REAL};
\node (fake)    [resultfake, below right=1.2cm and 1.0cm of dec] {FAKE};

\draw[arrow,color=primaryblue] (audio) -- (vad);
\draw[arrow,color=primaryblue] (vad)   -- (feat);
\draw[arrow,color=primaryblue] (feat)  -- (norm);
\draw[arrow,color=primaryblue] (norm)  -- (inf);
\draw[arrow,color=primaryblue] (inf)   -- (score);
\draw[arrow,color=primaryblue] (score) -- (dec);
\draw[arrow,color=primaryblue] (dec)   --
  node[above left,scale=0.8,color=successgreen]{Sim} (real);
\draw[arrow,color=primaryblue] (dec)   --
  node[above right,scale=0.8,color=dangerred]{N\~{a}o} (fake);

% === FASE TREINAMENTO (direita) ===
\draw[purpleaccent, line width=1.5pt, rounded corners=6pt, fill=purpleaccent!4]
  (0.5,6.8) rectangle (6.8,-4.0);
\node[font=\footnotesize\bfseries, color=purpleaccent] at (3.65,6.3)
  {TREINAMENTO};

\node (dataset) [startend]  at (3.65, 5.5) {Dataset};
\node (prep)    [block, below=of dataset]  {Pr\'{e}-processamento};
\node (ftrain)  [block, below=of prep]     {Extra\c{c}\~{a}o de Features};
\node (arch)    [block, below=of ftrain]   {Arquiteturas Neurais};
\node (train)   [block, below=of arch]     {Treinamento};
\node (conv)    [decision, below=of train] {Convergiu?};
\node (saved)   [modelnode, below=of conv] {Modelo Salvo};

\draw[arrow,color=purpleaccent] (dataset) -- (prep);
\draw[arrow,color=purpleaccent] (prep)    -- (ftrain);
\draw[arrow,color=purpleaccent] (ftrain)  -- (arch);
\draw[arrow,color=purpleaccent] (arch)    -- (train);
\draw[arrow,color=purpleaccent] (train)   -- (conv);
\draw[arrow,color=purpleaccent] (conv)    --
  node[right,scale=0.8,color=successgreen]{Sim} (saved);
\draw[arrow,color=purpleaccent] (conv.east) -| ++(1.2,0) |- (train.east)
  node[pos=0.25, right, scale=0.8, color=dangerred]{N\~{a}o};

% === CONEXAO DEPLOY ===
\draw[arrow, dashed, line width=1.2pt, color=purpleaccent]
  (saved.south west) .. controls +(0,-1.0) and +(2,1.0) .. (inf.east)
  node[pos=0.45, below, scale=0.8, color=purpleaccent]{Deploy};

\end{tikzpicture}
\caption{Diagrama do pipeline proposto, com fases de treinamento e
predi\c{c}\~{a}o.}
\label{fig:pipeline}
\end{figure}

\subsection{Etapa 1: Pr\'{e}-processamento}

O sistema suporta m\'{u}ltiplos formatos (WAV, MP3, FLAC, OGG) com
convers\~{a}o autom\'{a}tica para padr\~{o}es consistentes: taxa de amostragem
de 16\,kHz, mono, 16\,bits PCM linear.

\subsubsection{Normaliza\c{c}\~{a}o AGC}

A normaliza\c{c}\~{a}o implementa controle autom\'{a}tico de ganho adaptativo:

\begin{equation}
x_{norm}[n] = x[n] \cdot
\frac{L_{target}}{\sqrt{\dfrac{1}{N}\displaystyle\sum_{i=0}^{N-1}x[i]^2}}
\end{equation}

onde $L_{target} = -23$\,LUFS \'{e} o n\'{i}vel alvo de \textit{loudness}.

\subsubsection{Detec\c{c}\~{a}o de Atividade de Voz (VAD)}

O VAD baseado em energia remove segmentos silenciosos:

\begin{equation}
E[m] = 10\log_{10}\!\left(\frac{1}{N}\sum_{n=0}^{N-1}|x[mH+n]|^2\right)
\end{equation}

\begin{equation}
\text{Voice}[m] = \begin{cases}
1, & \text{se } E[m] > \theta_{VAD} \text{ e } D[m] > D_{min}\\
0, & \text{caso contr\'{a}rio}
\end{cases}
\end{equation}

onde $\theta_{VAD} = -25$\,dB, $H$ \'{e} o \textit{hop length}, $D[m]$ \'{e}
a dura\c{c}\~{a}o do segmento e $D_{min} = 0{,}3$\,s.

\begin{infobox}{VAD em Ambiente de Produ\c{c}\~{a}o}
O Silero VAD (disponibilizado via Torch Hub) oferece performance pr\'{e}-treinada
superior ao VAD por energia. Nos experimentos, a indisponibilidade do
\texttt{torchaudio} motivou o uso do VAD por energia como \textit{fallback},
preservando a funcionalidade do pipeline sem depend\^{e}ncias externas
n\~{a}o-atendidas no ambiente de CPU.
\end{infobox}

\subsection{Etapa 2: Extra\c{c}\~{a}o e An\'{a}lise de Caracter\'{i}sticas}

O sistema extrai caracter\'{i}sticas ac\'{u}sticas complementares cuja
hierarquia discriminativa \'{e} suportada pela literatura consolidada de
\textit{anti-spoofing} \cite{sahidullah2015comparison, todisco2016new,
todisco2019asvspoof}.

\subsubsection{Caracter\'{i}sticas Espectrais}

\textbf{Centroide Espectral}:
\begin{equation}
SC[m] = \frac{\displaystyle\sum_{k=0}^{K-1} k\cdot|X[m,k]|^2}
             {\displaystyle\sum_{k=0}^{K-1}|X[m,k]|^2}
\end{equation}

\textbf{Rolloff Espectral}: frequ\^{e}ncia abaixo da qual est\'{a} 85\% da
energia
\begin{equation}
SR[m] = k_r \text{ tal que }
\sum_{k=0}^{k_r}|X[m,k]|^2 = 0{,}85\sum_{k=0}^{K-1}|X[m,k]|^2
\end{equation}

\textbf{Largura de Banda Espectral}:
\begin{equation}
SB[m] =
\sqrt{\frac{\displaystyle\sum_{k=0}^{K-1}(k-SC[m])^2\cdot|X[m,k]|^2}
           {\displaystyle\sum_{k=0}^{K-1}|X[m,k]|^2}}
\end{equation}

\textbf{Taxa de Cruzamento Zero (ZCR)}:
\begin{equation}
ZCR[m] = \frac{1}{2N}\sum_{n=0}^{N-1}
|\mathrm{sgn}(x[mH+n]) - \mathrm{sgn}(x[mH+n-1])|
\end{equation}

\textbf{Contraste Espectral}:
\begin{equation}
SContrast[m] = \log\!\left(\frac{\mu_p}{\mu_v}\right)
\end{equation}

onde $\mu_p$ e $\mu_v$ s\~{a}o as m\'{e}dias dos picos e vales espectrais.

\textbf{Planura Espectral}:
\begin{equation}
SFlatness[m] =
\frac{\left(\displaystyle\prod_{k=0}^{K-1}|X[m,k]|^2\right)^{1/K}}
     {\dfrac{1}{K}\displaystyle\sum_{k=0}^{K-1}|X[m,k]|^2}
\end{equation}

\subsubsection{Mel-Frequency Cepstral Coefficients (MFCC)}

\begin{equation}
MFCC[m,c] = \sum_{k=0}^{M-1}\log(S[m,k])\cdot
\cos\!\left(\frac{\pi c(k+0{,}5)}{M}\right)
\end{equation}

onde $S[m,k]$ s\~{a}o as energias dos filtros mel, $M=40$ filtros,
coeficientes: 40, pr\'{e}-\^{e}nfase $\alpha=0{,}97$, faixa
80\,Hz--8\,kHz.

\subsubsection{Mel Spectrograms}

\begin{equation}
MelSpec[m,k] = \log\!\left(\sum_{n}H_k[n]\cdot|STFT[m,n]|^2 +
\varepsilon\right)
\end{equation}

Configura\c{c}\~{o}es: 80 bandas mel, janela Hamming de 1024 amostras,
\textit{hop length} 256, N\_FFT 2048.

\subsubsection{Caracter\'{i}sticas Pros\'{o}dicas}

\textbf{Frequ\^{e}ncia Fundamental (F0)}: extra\'{i}da por algoritmo CREPE
baseado em correla\c{c}\~{a}o de harm\^{o}nicos.

\textbf{Jitter} (varia\c{c}\~{a}o de per\'{i}odo ciclo a ciclo):
\begin{equation}
Jitter = \frac{1}{N-1}\sum_{i=1}^{N-1}\frac{|T_i - T_{i-1}|}{\bar{T}}
\end{equation}

\textbf{Shimmer} (varia\c{c}\~{a}o de amplitude ciclo a ciclo):
\begin{equation}
Shimmer = \frac{1}{N-1}\sum_{i=1}^{N-1}\frac{|A_i - A_{i-1}|}{\bar{A}}
\end{equation}

\textbf{HNR} (\textit{Harmonic-to-Noise Ratio}):
\begin{equation}
HNR = 10\log_{10}\!\left(\frac{P_{harmonic}}{P_{noise}}\right)
\end{equation}

\subsubsection{Constant-Q Transform (CQT)}

\begin{equation}
CQT[m,k] = \sum_{n}x[n]\cdot g_k[n-mH]\cdot e^{-j2\pi Q_k n/N_k}
\end{equation}

O CQCC derivado do CQT foi estabelecido como \textit{front-end} de
refer\^{e}ncia para \textit{anti-spoofing} por Todisco et al.\
\cite{todisco2016new}.
Configura\c{c}\~{o}es: 12 \textit{bins} por oitava, 7 oitavas
(84 caracter\'{i}sticas), $f_{min}=65$\,Hz.

\subsubsection{Caracter\'{i}sticas Delta e Delta-Delta}

\begin{equation}
\Delta X[m,c] = \frac{\displaystyle\sum_{n=1}^{N}n(X[m+n,c]-X[m-n,c])}
                     {2\displaystyle\sum_{n=1}^{N}n^2}
\end{equation}

\subsection{Ranking de Caracter\'{i}sticas por Desempenho Discriminativo}

A escolha das caracter\'{i}sticas ac\'{u}sticas \'{e} ancorada na literatura
de \textit{anti-spoofing}.
Caracter\'{i}sticas espectrais apresentam consistentemente maior poder
discriminativo em rela\c{c}\~{a}o a MFCCs cl\'{a}ssicos, conforme demonstrado
em estudos sistem\'{a}ticos conduzidos nas edi\c{c}\~{o}es do ASVspoof
\cite{sahidullah2015comparison, todisco2016new, todisco2019asvspoof}.
Caracter\'{i}sticas pros\'{o}dicas, embora complementares, exibem desempenho
inferior quando utilizadas isoladamente, resultado consistente com os achados
do ADD Challenge \cite{yi2022add}.

\begin{table}[H]
\centering
\caption{Ranking de Caracter\'{i}sticas Ac\'{u}sticas: S\'{i}ntese da
Literatura}
\label{tab:ranking_features}
\small
\begin{tabular}{p{2.8cm}ccp{4.8cm}}
\toprule
\textbf{Tipo} & \textbf{Dimen.} & \textbf{Custo} &
\textbf{Valida\c{c}\~{a}o na Literatura} \\
\midrule
Espectrais
  & 14 & M\'{e}dio
  & Maior poder discriminativo em detec\c{c}\~{a}o de artefatos de
    \textit{vocoder}~\cite{bartusiak2022, todisco2016new} \\
\addlinespace
Combina\c{c}\~{a}o otimizada
  & 60 & Alto
  & Fus\~{a}o de dom\'{i}nios reduz EER em cen\'{a}rios
    \textit{cross-dataset}~\cite{wang2024asvspoof5} \\
\addlinespace
MFCC
  & 40 & Baixo
  & Refer\^{e}ncia cl\'{a}ssica; LFCC supera MFCC em \textit{spoofing}
    desde 2015~\cite{sahidullah2015comparison, todisco2019asvspoof} \\
\addlinespace
Mel Spectrogram
  & 80 & Baixo
  & Base para CNNs; captura artefatos visuais de alta
    frequ\^{e}ncia~\cite{jung2022aasist, tak2021rawnet2} \\
\addlinespace
CQT
  & 84 & Alto
  & Resolu\c{c}\~{a}o logar\'{i}tmica; CQCC como \textit{front-end}
    robusto para \textit{anti-spoofing}~\cite{todisco2016new} \\
\addlinespace
Pros\'{o}dicas
  & 6 & M\'{e}dio
  & Desempenho isolado inferior; fus\~{a}o com espectrais melhora
    resultados~\cite{yi2022add, muller2022warning} \\
\bottomrule
\end{tabular}
\end{table}

\subsection{Normaliza\c{c}\~{a}o Estat\'{i}stica}

\textbf{Min-Max Scaling}:
\begin{equation}
X_{norm} = \frac{X - X_{min}}{X_{max} - X_{min}}
\end{equation}

\textbf{Z-Score}:
\begin{equation}
X_{norm} = \frac{X - \mu}{\sigma}
\end{equation}

% =====================================================================
%  4. ARQUITETURAS DE DEEP LEARNING
% =====================================================================
\section{Arquiteturas de \textit{Deep Learning}}
\label{sec:arquiteturas}

O pipeline implementa cinco fam\'{i}lias de arquiteturas neurais, cada uma
otimizada para capturar aspectos distintos dos padr\~{o}es de \'{a}udio
sint\'{e}tico.
As arquiteturas foram implementadas em TensorFlow~2.21.0 e Keras~3 com
interfaces unificadas de treinamento e infer\^{e}ncia.

\subsection{Fam\'{i}lia 1: Modelos Foundation Pr\'{e}-treinados (SSL)}

\subsubsection{WavLM e HuBERT}

O WavLM \cite{chen2022wavlm} e o HuBERT \cite{hsu2021hubert} s\~{a}o
pr\'{e}-treinados em larga escala e especializados para tarefas de voz via
\textit{fine-tuning}.
Neste pipeline, ambos s\~{a}o suportados como cabe\c{c}alho de
classifica\c{c}\~{a}o sobre representa\c{c}\~{o}es SSL congeladas.
A integra\c{c}\~{a}o \'{e} feita com implementa\c{c}\~{o}es simplificadas
compat\'{i}veis com o ambiente sem GPU, exigindo GPU para o desempenho
reportado na literatura \cite{wang2024asvspoof5}.

\subsection{Fam\'{i}lia 2: An\'{a}lise Direta de Forma de Onda e Grafos}

\subsubsection{AASIST}

O AASIST \cite{jung2022aasist} implementa Redes de Aten\c{c}\~{a}o em Grafos
Heterog\^{e}neos (HAGs) para modelar depend\^{e}ncias espectrais e temporais
simultaneamente.
A implementa\c{c}\~{a}o nativa inclui SincConv e camadas de grafos com
1.150.592 par\^{a}metros.
\textbf{Resultado experimental}: n\~{a}o convergiu em ambiente CPU com
420 amostras de treino (acur\'{a}cia de valida\c{c}\~{a}o estabilizada em
50\%).

\subsubsection{RawNet2}

O RawNet2 \cite{tak2021rawnet2} processa diretamente a forma de onda PCM 1D
via \textit{Sinc-Convolutions}, com 5.303.553 par\^{a}metros.
\textbf{Resultado experimental}: n\~{a}o convergiu, com AUC-ROC de 0,000 no
conjunto de teste, indicando colapso do classificador em ambiente CPU.

\subsection{Fam\'{i}lia 3: EfficientNet-LSTM}

Combina EfficientNet para extra\c{c}\~{a}o de caracter\'{i}sticas espectrais
com LSTM para modelagem temporal.
O espectrograma mel \'{e} redimensionado para $224\times224$ e replicado em 3
canais para compatibilidade com EfficientNet.
A aten\c{c}\~{a}o temporal \'{e} calculada como:

\begin{equation}
\alpha_t = \frac{\exp(e_t)}{\displaystyle\sum_{i=1}^{T}\exp(e_i)},
\qquad c = \sum_{t=1}^{T}\alpha_t h_t
\end{equation}

Com 7.986.852 par\^{a}metros e 30,5\,MB de mem\'{o}ria, \'{e} a maior
arquitetura em n\'{u}mero de par\^{a}metros.
\textbf{Resultado experimental}: n\~{a}o convergiu em ambiente CPU com
420 amostras (acur\'{a}cia de valida\c{c}\~{a}o estabilizada em 50\%).

\subsection{Fam\'{i}lia 4: MultiscaleCNN (Res2Net)}

A MultiscaleCNN utiliza como \textit{backbone} a Res2Net \cite{gao2021res2net},
que implementa convolu\c{c}\~{o}es paralelas em m\'{u}ltiplas escalas:

\begin{equation}
\text{MSConv}(X) =
\text{Concat}[\text{Conv}_{k_1}(X), \text{Conv}_{k_2}(X),
\text{Conv}_{k_3}(X)]
\end{equation}

onde $k_1=3$, $k_2=5$, $k_3=7$.
A configura\c{c}\~{a}o Res2Net-50 (escala=8, baseWidth=26) totaliza
48.965.581 par\^{a}metros e 186,8\,MB de mem\'{o}ria.
\textbf{Resultado experimental}: melhor modelo individual --- acur\'{a}cia de
97,78\%, EER de 5,56\%, AUC-ROC de 0,997, lat\^{e}ncia de 133,7\,ms.

\subsection{Fam\'{i}lia 5: Ensemble Adaptativo}

O Ensemble Adaptativo implementa 5 ramos convolucionais paralelos com
fus\~{a}o ponderada:

\begin{equation}
w_i = \mathrm{softmax}\!\left(\text{Dense}\!\left(
\mathrm{concat}[f_1,\ldots,f_N]\right)\right)_i
\end{equation}

\begin{equation}
\text{Ensemble}(x) = \sum_{i=1}^{N}w_i\cdot\text{Model}_i(x)
\end{equation}

Com apenas 2.340.650 par\^{a}metros e 8,9\,MB, \'{e} a arquitetura mais
leve que convergiu.
\textbf{Resultado experimental}: acur\'{a}cia de 91,11\%, EER de 13,33\%,
AUC-ROC de 0,975, lat\^{e}ncia de 48,1\,ms.

% =====================================================================
%  5. PIPELINE DE CODIGO ABERTO
% =====================================================================
\section{Proposta de Pipeline de C\'{o}digo Aberto}
\label{sec:pipeline_opensource}

O pipeline foi projetado para ser plenamente reprodut\'{i}vel com ferramentas
livres.
A Figura~\ref{fig:pipeline_compact} mapeia os componentes recomendados.

\begin{figure}[H]
\centering
\begin{tikzpicture}[
  node distance=1cm,
  every node/.style={font=\scriptsize},
  block/.style={
    rectangle, rounded corners=2pt,
    draw=blue!60, fill=blue!5,
    minimum width=3.5cm, minimum height=0.7cm,
    text centered
  },
  arrow/.style={->,>=stealth,thick}
]

\node (in)   [block] {\'{A}udio (WAV/MP3/FLAC)};
\node (qc)   [block, below of=in] {Valida\c{c}\~{a}o (SNR / Clipping)};
\node (prep) [block, below of=qc] {Pr\'{e}-proc.\ (VAD + AGC + Resample)};

\node (feat) at (-4,-4) [block] {Features \\ MFCC / CQT};
\node (ssl)  at (0,-4)  [block] {SSL \\ WavLM / HuBERT};
\node (raw)  at (4,-4)  [block] {Waveform \\ Raw};

\node (m1) at (-4,-5.5) [block] {MLP / CNN};
\node (m2) at (0,-5.5)  [block] {AASIST / Conformer};
\node (m3) at (4,-5.5)  [block] {RawNet2};

\node (fusion) at (0,-7) [block] {Fusion \\ XGBoost / Attention};

\node (post) at (0,-8.2) [block] {Smoothing Temporal};
\node (out)  at (0,-9.4) [block] {REAL / FAKE + Confian\c{c}a};

\draw[arrow] (in) -- (qc);
\draw[arrow] (qc) -- (prep);
\draw[arrow] (prep) -- (feat);
\draw[arrow] (prep) -- (ssl);
\draw[arrow] (prep) -- (raw);
\draw[arrow] (feat) -- (m1);
\draw[arrow] (ssl) -- (m2);
\draw[arrow] (raw) -- (m3);
\draw[arrow] (m1) -- (fusion);
\draw[arrow] (m2) -- (fusion);
\draw[arrow] (m3) -- (fusion);
\draw[arrow] (fusion) -- (post);
\draw[arrow] (post) -- (out);

\end{tikzpicture}
\caption{Pipeline compacto para detec\c{c}\~{a}o de \textit{deepfake} de voz.}
\label{fig:pipeline_compact}
\end{figure}

\subsection{Componentes de Software Livre}

\begin{itemize}
    \item \textbf{Pr\'{e}-processamento e VAD:} Silero VAD (via Torch Hub) com
    fallback por energia para ambientes sem \texttt{torchaudio}. Para AGC e
    reamostragem, \texttt{librosa} prov\^{e} opera\c{c}\~{o}es eficientes em
    CPU.
    \item \textbf{Extra\c{c}\~{a}o de Caracter\'{i}sticas:} \texttt{librosa}
    e \texttt{essentia} para extra\c{c}\~{a}o de caracter\'{i}sticas espectrais,
    MFCC e pros\'{o}dicas.
    \item \textbf{Modelos Foundation:} Dispon\'{i}veis via
    \texttt{transformers} (Hugging Face). Exigem GPU para fine-tuning
    eficiente.
    \item \textbf{Arquiteturas Especializadas:} RawNet2 e AASIST com
    implementa\c{c}\~{o}es oficiais no GitHub sob licen\c{c}a MIT.
    \item \textbf{Ensemble e Classifica\c{c}\~{a}o:} TensorFlow/Keras para
    implementa\c{c}\~{a}o dos modelos e fus\~{a}o adaptativa.
    \item \textbf{Interpretabilidade (proposta):} \texttt{shap} para valores
    SHAP e \texttt{captum} (PyTorch) para GradCAM --- integra\c{c}\~{a}o
    planejada como extens\~{a}o futura.
\end{itemize}

\subsection{Configura\c{c}\~{a}o de Treinamento}

\subsubsection{Fun\c{c}\~{a}o de Perda}

\textit{Binary Cross-Entropy} com balanceamento de classes:

\begin{equation}
\mathcal{L} = -\frac{1}{N}\sum_{i=1}^{N}\left[
w_1 y_i\log(\hat{y}_i) + w_0(1-y_i)\log(1-\hat{y}_i)\right]
\end{equation}

\subsubsection{Regulariza\c{c}\~{a}o}

Regulariza\c{c}\~{a}o L2 com $\lambda=0{,}0001$, \textit{Early Stopping}
com paci\^{e}ncia de 10 \'{e}pocas, \textit{ReduceLROnPlateau} com
paci\^{e}ncia de 5 \'{e}pocas.

\subsubsection{Otimizador}

Adam com $\alpha=0{,}001$, $\beta_1=0{,}9$, $\beta_2=0{,}999$.

% =====================================================================
%  6. EXPERIMENTOS E RESULTADOS
% =====================================================================
\section{Experimentos e Resultados}
\label{sec:resultados}

\subsection{Dataset e Configura\c{c}\~{a}o Experimental}

\begin{warningbox}{Configura\c{c}\~{a}o Experimental Real}
Todos os resultados desta se\c{c}\~{a}o foram obtidos experimentalmente no
mesmo ambiente de execu\c{c}\~{a}o: CPU Windows 11 (sem GPU), Python~3.11,
TensorFlow~2.21.0, Keras~3. Os valores reproduzem com fidelidade as
m\'{e}tricas medidas durante o treinamento e os testes de robustez.
\end{warningbox}

O \textit{dataset} BRSpeech-DF foi constru\'{i}do com 600 amostras
balanceadas de 1\,segundo cada (16\,kHz, mono), correspondendo a fala
portuguesa real (\textit{bonafide}) e sint\'{e}tica (\textit{spoof}).
As amostras foram submetidas ao pr\'{e}-processamento com VAD por energia
(fallback do Silero VAD, indispon\'{i}vel sem \texttt{torchaudio}) e AGC
baseada em normaliza\c{c}\~{a}o RMS para $-23$\,LUFS.

\begin{itemize}
    \item \textbf{Distribui\c{c}\~{a}o}: 420 treino (210 real + 210 fake),
    90 valida\c{c}\~{a}o (45+45), 90 teste (45+45) --- split
    70/15/15 estratificado.
    \item \textbf{Hardware}: CPU (sem GPU --- Windows 11, AMD/Intel x86-64)
    \item \textbf{\textit{Framework}}: TensorFlow 2.21.0 + Keras 3,
    Python 3.11
    \item \textbf{Batch Size}: 32
    \item \textbf{\'{E}pocas m\'{a}x.}: 20 (\textit{early stopping}
    paci\^{e}ncia = 10)
    \item \textbf{\textit{Data Augmentation}}: \textit{pitch shifting}
    ($\pm$2 semitons), \textit{time stretching} ($0{,}9$--$1{,}1\times$)
\end{itemize}

\subsection{M\'{e}tricas de Avalia\c{c}\~{a}o}

\textbf{Equal Error Rate (EER)}:
\begin{equation}
EER = \tau \text{ tal que } FAR(\tau) = FRR(\tau)
\end{equation}

\textbf{AUC-ROC}:
\begin{equation}
AUC = \int_0^1 TPR(FPR^{-1}(x))\,dx
\end{equation}

EER foi calculado via \texttt{scipy.optimize.brentq} sobre a curva ROC;
AUC-ROC via \texttt{sklearn.metrics.roc\_auc\_score}.

\subsection{Resultados por Arquitetura}

A Tabela~\ref{tab:resultados} apresenta o desempenho quantitativo das cinco
arquiteturas implementadas, com resultados medidos no conjunto de teste
(90 amostras, 45 real + 45 fake).
Duas arquiteturas convergiram com sucesso: MultiscaleCNN e Ensemble
Adaptativo --- ambas baseadas em espectrograma.
Tr\^{e}s arquiteturas de an\'{a}lise direta de forma de onda
(EfficientNet-LSTM, AASIST, RawNet2) n\~{a}o convergiram, apresentando
acur\'{a}cia equivalente a acaso (50\%), resultado consistente com os
requisitos computacionais documentados na literatura para esse tipo de
modelo \cite{tak2021rawnet2, jung2022aasist}.

\begin{table}[H]
\centering
\caption{Desempenho das Arquiteturas Implementadas (conjunto de teste,
90 amostras)}
\label{tab:resultados}
\small
\begin{tabular}{lcccccc}
\toprule
\textbf{Arquitetura} & \textbf{Acur.} & \textbf{EER} &
\textbf{AUC-ROC} & \textbf{Lat.\,(ms)} & \textbf{\'{E}poc.} &
\textbf{Converge?} \\
\midrule
%% AUTOGEN_BEGIN:tab_resultados
\textbf{MultiscaleCNN (Res2Net-50)}
  & \textbf{97,78\%} & \textbf{5,56\%} & \textbf{0,997}
  & 133,7 & 20 & \textcolor{successgreen}{Sim} \\
\textbf{Ensemble Adaptativo}
  & \textbf{91,11\%} & \textbf{13,33\%} & \textbf{0,975}
  & 48,1 & 20 & \textcolor{successgreen}{Sim} \\
EfficientNet-LSTM
  & 50,00\% & 46,67\% & 0,511
  & 63,8 & 11 & \textcolor{dangerred}{N\~{a}o} \\
AASIST
  & 50,00\% & 50,00\% & 0,500
  & 89,0 & 11 & \textcolor{dangerred}{N\~{a}o} \\
RawNet2
  & 50,00\% & 50,00\% & 0,000
  & 156,4 & 11 & \textcolor{dangerred}{N\~{a}o} \\
%% AUTOGEN_END:tab_resultados
\bottomrule
\end{tabular}
\end{table}

\begin{infobox}{Interpreta\c{c}\~{a}o da N\~{a}o-Converg\^{e}ncia}
A estabiliza\c{c}\~{a}o em 50\% de acur\'{a}cia e a sa\'{i}da do
\textit{early stopping} na \'{e}poca~11 indicam colapso do gradiente em todos
os modelos \textit{raw-audio}.
O AUC-ROC de 0,000 do RawNet2 revela que as probabilidades preditas s\~{a}o
inversas \`{a} classe real --- o classificador aprendeu um padr\~{a}o
espúrio sem generaliza\c{c}\~{a}o.
A literatura confirma que RawNet2 e AASIST exigem GPU e $>$5.000 amostras:
sem esses recursos, a complexidade das SincConv e GATs n\~{a}o se justifica
para a capacidade de generaliza\c{c}\~{a}o dispon\'{i}vel.
\end{infobox}

\subsection{Efici\^{e}ncia Computacional}

A Tabela~\ref{tab:complexidade} compara par\^{a}metros, mem\'{o}ria e
lat\^{e}ncia medidas experimentalmente para \'{a}udios de 10\,segundos
em CPU.

\begin{table}[H]
\centering
\caption{An\'{a}lise de Complexidade Computacional (medido em CPU)}
\label{tab:complexidade}
\begin{tabular}{lcccc}
\toprule
\textbf{Arquitetura} & \textbf{Par\^{a}metros} &
\textbf{Mem\'{o}ria (MB)} & \textbf{Lat\^{e}ncia (ms)} &
\textbf{Converge (CPU)} \\
\midrule
%% AUTOGEN_BEGIN:tab_complexidade
MultiscaleCNN & 48.965.581 & 186,8 & 133,7 & \textcolor{successgreen}{Sim} \\
Effic.-LSTM & 7.986.852 & 30,5 & 63,8 & \textcolor{dangerred}{N\~{a}o} \\
RawNet2 & 5.303.553 & 20,2 & 156,4 & \textcolor{dangerred}{N\~{a}o} \\
Ensemble Adapt. & 2.340.650 & 8,9 & 48,1 & \textcolor{successgreen}{Sim} \\
AASIST & 1.150.592 & 4,4 & 89,0 & \textcolor{dangerred}{N\~{a}o} \\
%% AUTOGEN_END:tab_complexidade
\bottomrule
\end{tabular}
\end{table}

Nota-se que o Ensemble Adaptativo, com apenas 8,9\,MB e 48,1\,ms de
lat\^{e}ncia, \'{e} o modelo convergente mais eficiente --- adequado para
aplica\c{c}\~{o}es com restri\c{c}\~{o}es de hardware.
A MultiscaleCNN, embora mais pesada (186,8\,MB), oferece o melhor desempenho
absoluto entre os modelos que convergiram.

\subsection{Testes de Robustez sob Ru\'{i}do AWGN}

Para avaliar o comportamento das arquiteturas em condi\c{c}\~{o}es reais de
uso, os dois modelos convergentes foram submetidos a testes de robustez com
ru\'{i}do gaussiano branco aditivo (AWGN) em tr\^{e}s n\'{i}veis de
SNR: 10\,dB, 20\,dB e 30\,dB.
Os 90 exemplos do conjunto de teste foram corrompidos de forma sint\'{e}tica
antes da infer\^{e}ncia; os pesos dos modelos permaneceram congelados.

\begin{table}[H]
\centering
\caption{Tabela~5 --- Robustez sob Ru\'{i}do AWGN: Acur\'{a}cia (\%) e EER (\%)}
\label{tab:robustez}
\begin{tabular}{lcccccccc}
\toprule
\multirow{2}{*}{\textbf{Arquitetura}} &
\multicolumn{2}{c}{\textbf{Limpo}} &
\multicolumn{2}{c}{\textbf{SNR 30\,dB}} &
\multicolumn{2}{c}{\textbf{SNR 20\,dB}} &
\multicolumn{2}{c}{\textbf{SNR 10\,dB}} \\
\cmidrule(lr){2-3}\cmidrule(lr){4-5}\cmidrule(lr){6-7}\cmidrule(lr){8-9}
 & Acur. & EER & Acur. & EER & Acur. & EER & Acur. & EER \\
\midrule
%% AUTOGEN_BEGIN:tab_robustez
MultiscaleCNN
  & 97,78\% & 2,22\%
  & 71,11\% & 28,89\%
  & 72,22\% & 27,78\%
  & 58,89\% & 41,11\% \\
Ensemble Adapt.
  & 87,78\% & 12,22\%
  & 74,44\% & 25,56\%
  & 50,00\% & 50,00\%
  & 50,00\% & 50,00\% \\
%% AUTOGEN_END:tab_robustez
\bottomrule
\end{tabular}
\end{table}

Os resultados revelam padr\~{o}es distintos de degrada\c{c}\~{a}o:

\begin{itemize}
    \item \textbf{MultiscaleCNN}: degrada gradualmente e de forma n\~{a}o
    monotônica (72,22\% em SNR~20\,dB supera ligeiramente os 71,11\% em
    SNR~30\,dB), indicando que o modelo \'{e} mais robusto a n\'{i}veis
    moderados de ru\'{i}do. A 10\,dB, a acur\'{a}cia cai para 58,89\%,
    ainda acima do acaso.
    \item \textbf{Ensemble Adaptativo}: mantém desempenho razo\'{a}vel em
    SNR~30\,dB (74,44\%), mas colapsa abaixo de SNR~20\,dB
    (50,00\%), indicando alta sensibilidade ao ru\'{i}do.
    O Ensemble, por combinar m\'{u}ltiplos ramos com pesos adaptativos,
    perde o sinal discriminativo quando ru\'{i}do cont\'{a}mina todos os
    ramos simultaneamente.
\end{itemize}

\begin{warningbox}{Implica\c{c}\~{a}o para Uso em Produ\c{c}\~{a}o}
Nenhum dos modelos testados manteve desempenho aceit\'{a}vel em SNR~10\,dB
(EER $>$ 40\%). Para uso em ambientes ruidosos (e.g., telefonas, grava\c{c}\~{o}es
ambiente), \'{e} indispens\'{a}vel \textit{data augmentation} com ru\'{i}do
durante o treinamento ou a integra\c{c}\~{a}o de um m\'{o}dulo de
melhoria de fala (\textit{speech enhancement}) como pr\'{e}-est\'{a}gio.
\end{warningbox}

\subsection{Interpretabilidade por Valores SHAP (Extens\~{a}o Proposta)}

A interpretabilidade dos modelos por valores SHAP (\textit{SHapley Additive
exPlanations}) \cite{lundberg2017unified}, embora n\~{a}o implementada
experimentalmente neste trabalho, constitui uma extens\~{a}o natural e
priorit\'{a}ria do pipeline proposto.

O m\'{e}todo SHAP fundamenta-se na \textbf{Teoria dos Jogos Cooperativos}:
trata a predi\c{c}\~{a}o como o resultado de um jogo, onde cada
caracter\'{i}stica ac\'{u}stica \'{e} um ``jogador'' e o valor SHAP
$\phi_i$ representa sua contribui\c{c}\~{a}o marginal justa:

\begin{equation}
\phi_i = \sum_{S \subseteq F \setminus \{i\}}
\frac{|S|!\,(|F|-|S|-1)!}{|F|!}
\left[f(S \cup \{i\}) - f(S)\right]
\end{equation}

Com base nos achados da literatura \cite{bartusiak2022, todisco2016new},
espera-se que a aplica\c{c}\~{a}o do SHAP sobre a MultiscaleCNN confirme a
predomin\^{a}ncia das caracter\'{i}sticas espectrais de alta frequ\^{e}ncia
--- regi\~{a}o onde vocoders neurais deixam artefatos residuais
subliminares --- sobre pistas pros\'{o}dicas, que sistemas TTS modernos
replicam com alta fidelidade.

A integra\c{c}\~{a}o da an\'{a}lise SHAP ao pipeline est\'{a} descrita
na Se\c{c}\~{a}o de Trabalhos Futuros.

\subsection{Compara\c{c}\~{a}o com M\'{e}todos da Literatura}

\begin{warningbox}{Nota Metodol\'{o}gica}
Os resultados do sistema proposto foram obtidos em \textit{dataset} pr\'{o}prio
(600 amostras, BRSpeech-DF), em ambiente CPU, enquanto os m\'{e}todos da
literatura foram avaliados nos \textit{benchmarks} ASVspoof 2019/2021 com GPU.
A compara\c{c}\~{a}o direta \'{e} metodologicamente inadequada.
A tabela objetiva contextualizar o posicionamento relativo das arquiteturas
adotadas em rela\c{c}\~{a}o ao estado da arte.
\end{warningbox}

\begin{table}[H]
\centering
\caption{Arquiteturas adotadas e desempenho reportado em
\textit{benchmarks} p\'{u}blicos}
\label{tab:comparacao}
\small
\begin{tabular}{p{3.5cm}ccp{3.5cm}c}
\toprule
\textbf{M\'{e}todo} & \textbf{EER (\%)} & \textbf{AUC} &
\textbf{Dataset} & \textbf{Refer\^{e}ncia} \\
\midrule
GMM-UBM + CQCC   & 22,1  & --- & ASVspoof 2019 LA & \cite{todisco2019asvspoof} \\
LCNN + MFCC      & 14,6  & --- & ASVspoof 2021 LA & \cite{yamagishi2021asvspoof} \\
RawNet2          & 7,6   & --- & ASVspoof 2021 LA & \cite{tak2021rawnet2} \\
AASIST           & 0,83  & --- & ASVspoof 2019 LA & \cite{jung2022aasist} \\
WavLM + MLP      & 3,12  & --- & ASVspoof 5       & \cite{wang2024asvspoof5} \\
\midrule
%% AUTOGEN_BEGIN:tab_comparacao_our
\textbf{MultiscaleCNN (este trab.)}
  & \textbf{5,56} & \textbf{0,997}
  & BRSpeech-DF (CPU) & --- \\
\textbf{Ensemble Adapt. (este trab.)}
  & \textbf{13,33} & \textbf{0,975}
  & BRSpeech-DF (CPU) & --- \\
%% AUTOGEN_END:tab_comparacao_our
\bottomrule
\end{tabular}
\end{table}

A MultiscaleCNN alcan\c{c}a EER de 5,56\% em ambiente CPU com \textit{dataset}
reduzido, posicionando-se entre RawNet2 (7,6\%, GPU + ASVspoof) e AASIST
(0,83\%, GPU + ASVspoof) nas arquiteturas da 3\textordfeminine\ gera\c{c}\~{a}o
--- resultado not\'{a}vel considerando as restri\c{c}\~{o}es experimentais.

% =====================================================================
%  7. DISCUSSAO
% =====================================================================
\section{Discuss\~{a}o}
\label{sec:discussao}

\subsection{Dicotomia Espectrograma vs.\ Forma de Onda Bruta}

O principal achado emp\'{i}rico deste trabalho \'{e} a confirma\c{c}\~{a}o,
em ambiente controlado e reproducível, da dicotomia entre arquiteturas
baseadas em espectrograma e aquelas de an\'{a}lise direta de forma de onda:

\begin{itemize}
    \item \textbf{Modelos espectrogram-based} (MultiscaleCNN, Ensemble)
    convergiram com 420 amostras em CPU, atingindo acur\'{a}cia de
    97,78\% e 91,11\%, respectivamente.
    \item \textbf{Modelos \textit{raw-audio}} (EfficientNet-LSTM, AASIST,
    RawNet2) n\~{a}o convergiram nas mesmas condi\c{c}\~{o}es,
    estabilizando em 50\% de acur\'{a}cia --- performance de acaso.
\end{itemize}

Este resultado \'{e} consistente com o racioc\'{i}nio da segunda gera\c{c}\~{a}o
de m\'{e}todos: espectrogramas mel constituem representa\c{c}\~{o}es compactas
que amplificam visualmente os artefatos de alta frequ\^{e}ncia introduzidos
por vocoders, reduzindo a carga de aprendizado para a rede \cite{bartusiak2022}.
Modelos \textit{raw-audio}, ao trabalhar com PCM de 16.000 pontos por
segundo, precisam aprender essa transforma\c{c}\~{a}o implicitamente ---
tarefa que requer mais dados e maior capacidade computacional
\cite{tak2021rawnet2, jung2022aasist}.

\subsection{MultiscaleCNN: Melhor Modelo Individual}

A MultiscaleCNN (Res2Net-50) alcan\c{c}ou 97,78\% de acur\'{a}cia e
AUC-ROC de 0,997, superando o Ensemble Adaptativo em acur\'{a}cia mas
com maior lat\^{e}ncia (133,7\,ms vs.\ 48,1\,ms) e consumo de mem\'{o}ria
(186,8\,MB vs.\ 8,9\,MB).
A superioridade da Res2Net confirma a efic\'{a}cia das convolu\c{c}\~{o}es
multi-escala para capturar padr\~{o}es espectrais em m\'{u}ltiplas
granularidades \cite{gao2021res2net}.

\subsection{Ensemble: Melhor Rela\c{c}\~{a}o Efici\^{e}ncia-Desempenho}

O Ensemble Adaptativo, apesar de apresentar acur\'{a}cia inferior
(91,11\%) e EER mais alto (13,33\%), destaca-se por:
\begin{itemize}
    \item Menor lat\^{e}ncia: 48,1\,ms (adequado para aplica\c{c}\~{o}es em
    tempo quase real).
    \item Menor mem\'{o}ria: 8,9\,MB (implementa\c{c}\~{a}o embarcada
    vi\'{a}vel).
    \item Melhor robustez em SNR~30\,dB (74,44\%) em rela\c{c}\~{a}o
    \`{a} MultiscaleCNN (71,11\%).
\end{itemize}

\subsection{Degrada\c{c}\~{a}o sob Ru\'{i}do e Implica\c{c}\~{o}es}

A degrada\c{c}\~{a}o sob ru\'{i}do \'{e} substancial para ambos os modelos,
especialmente o Ensemble, que colapsa a 50\% para SNR $\leq$ 20\,dB.
Dois mecanismos explicam essa degrada\c{c}\~{a}o:

\begin{enumerate}
    \item \textbf{Sensibilidade de caracter\'{i}sticas espectrais ao
    ru\'{i}do}: centroide espectral, planura e contraste s\~{a}o fortemente
    perturbados por AWGN, pois o ru\'{i}do eleva uniformemente o piso
    espectral.
    \item \textbf{Aus\^{e}ncia de augmentation com ru\'{i}do}: os modelos
    foram treinados exclusivamente com augmentation de \textit{pitch} e
    \textit{time stretch}; a introdu\c{c}\~{a}o de ru\'{i}do durante o
    treinamento seria essencial para robustez.
\end{enumerate}

\subsection{Limitações}

\begin{itemize}
    \item \textbf{Tamanho do \textit{Dataset}}: 600 amostras (420 de
    treino) s\~{a}o insuficientes para modelos \textit{raw-audio};
    para modelos espectrogram-based, os resultados s\~{a}o promissores
    mas precisam ser validados em \textit{benchmarks} p\'{u}blicos.
    \item \textbf{Dataset monomodal}: treinamento exclusivamente com
    fala sintetizada via RVC~v2; generaliza\c{c}\~{a}o a outros geradores
    \'{e} desconhecida.
    \item \textbf{Aus\^{e}ncia de GPU}: modelos da 3\textordfeminine\
    e 4\textordfeminine\ gera\c{c}\~{a}o n\~{a}o podem ser avaliados
    adequadamente.
    \item \textbf{Robustez limitada}: nenhum modelo manteve EER aceit\'{a}vel
    em SNR~10\,dB.
    \item \textbf{SHAP n\~{a}o implementado}: interpreta\c{c}\~{a}o
    experimental das decis\~{o}es permanece como trabalho futuro.
\end{itemize}

\subsection{Trabalhos Futuros}

\begin{itemize}
    \item Avalia\c{c}\~{a}o nos \textit{benchmarks} p\'{u}blicos ASVspoof
    2019/2021 e ASVspoof~5 com GPU.
    \item \textit{Data augmentation} com ru\'{i}do AWGN e c\'{o}decagem MP3
    para melhoria de robustez.
    \item Re-treinamento com \textit{dataset} expandido ($>$5.000 amostras)
    para modelos \textit{raw-audio}.
    \item Integra\c{c}\~{a}o do an\'{a}lise SHAP para interpretabilidade
    local e global.
    \item Detec\c{c}\~{a}o de \textit{deepfakes} parciais com
    segmenta\c{c}\~{a}o temporal.
    \item Integra\c{c}\~{a}o de modelos multilinguais (mHuBERT).
\end{itemize}

% =====================================================================
%  8. CONCLUSOES
% =====================================================================
\section{Conclus\~{o}es}
\label{sec:conclusao}

Este trabalho desenvolveu e avaliou empiricamente um \textit{pipeline} modular
de c\'{o}digo aberto para detec\c{c}\~{a}o de \'{a}udio sint\'{e}tico,
conduzindo experimentos reais com cinco arquiteturas neurais sobre um
\textit{dataset} brasileiro balanceado de 600 amostras.

O resultado mais not\'{a}vel \'{e} a MultiscaleCNN (Res2Net-50), que atingiu
97,78\% de acur\'{a}cia e AUC-ROC de 0,997 em ambiente CPU sem GPU ---
desempenho competitivo mesmo diante de restri\c{c}\~{o}es severas de hardware.
O Ensemble Adaptativo complementa com 48,1\,ms de lat\^{e}ncia e
apenas 8,9\,MB de mem\'{o}ria, sendo vi\'{a}vel em ambientes embarcados.

A n\~{a}o-converg\^{e}ncia dos modelos \textit{raw-audio} (AASIST, RawNet2,
EfficientNet-LSTM) em CPU confirma empiricamente uma distin\c{c}\~{a}o
fundamental documentada na literatura: arquiteturas da
3\textordfeminine/4\textordfeminine\ gera\c{c}\~{a}o exigem GPU e
\textit{datasets} maiores para generalizar, enquanto modelos
espectrogram-based convergem com recursos modestos.
Este achado \'{e} relevante para praticantes que precisam escolher
arquiteturas em fun\c{c}\~{a}o de restri\c{c}\~{o}es de infraestrutura.

Os testes de robustez evidenciaram degrada\c{c}\~{a}o substancial sob
ru\'{i}do AWGN, estabelecendo a incorpora\c{c}\~{a}o de \textit{data
augmentation} com ru\'{i}do como prioridade imediata para uso em produ\c{c}\~{a}o.

As contribui\c{c}\~{o}es principais deste trabalho s\~{a}o:

\begin{enumerate}
    \item \textit{Survey} sistem\'{a}tico das quatro gera\c{c}\~{o}es de
    m\'{e}todos de detec\c{c}\~{a}o de \'{a}udio sint\'{e}tico.
    \item Ranking e valida\c{c}\~{a}o de caracter\'{i}sticas ac\'{u}sticas
    por capacidade discriminativa, ancorado na literatura.
    \item Desenvolvimento de um \textit{pipeline} modular de c\'{o}digo aberto
    com cinco arquiteturas neurais em TensorFlow/Keras.
    \item Avalia\c{c}\~{a}o emp\'{i}rica real com dicotomia clara entre
    modelos espectrogram-based (convergentes) e \textit{raw-audio}
    (n\~{a}o-convergentes em CPU).
    \item Testes de robustez sob ru\'{i}do AWGN em SNR~10, 20 e 30\,dB para
    os modelos convergentes.
    \item Proposta de extens\~{a}o via an\'{a}lise SHAP para
    interpretabilidade das decis\~{o}es.
\end{enumerate}

% =====================================================================
%  REFERENCIAS
% =====================================================================
\begin{thebibliography}{99}

\bibitem{baevski2020wav2vec}
BAEVSKI, A. et al.
wav2vec 2.0: A Framework for Self-Supervised Learning of Speech
Representations.
\textbf{Advances in Neural Information Processing Systems (NeurIPS)},
v.~33, p.~12449--12460, 2020.

\bibitem{bartusiak2022}
BARTUSIAK, E. R. et al.
Frequency Domain-based Detection of Generated Speech.
In: \textbf{IEEE International Workshop on Information Forensics and Security
(WIFS)}, p.~1--6, 2022.

\bibitem{bird2023realtime}
BIRD, J. J.; LOTFI, A.
Real-time Detection of AI-Generated Speech for DeepFake Voice Conversion.
\textbf{arXiv preprint} arXiv:2308.12734, 2023.
Dispon\'{i}vel em: \url{https://arxiv.org/abs/2308.12734}.
Acesso em: 15 jun.\ 2025.

\bibitem{chen2022wavlm}
CHEN, S. et al.
WavLM: Large-Scale Self-Supervised Pre-Training for Full Stack Speech
Processing.
\textbf{IEEE Journal of Selected Topics in Signal Processing},
v.~16, n.~6, p.~1505--1518, 2022.

\bibitem{dfrlab2024}
DFRLAB.
Brazil's electoral deepfake law tested as AI-generated content targeted local
elections.
\textbf{Digital Forensic Research Lab (Atlantic Council)}, 2024.
Dispon\'{i}vel em: \url{https://dfrlab.org/2024/03/brazil-deepfake-elections}.
Acesso em: 15 jun.\ 2025.

\bibitem{farrugia2024}
FARRUGIA, B.
The challenges of identifying deepfakes ahead of the 2024 Brazil election.
\textbf{Digital Forensic Research Lab (Atlantic Council)}, 2024.
Dispon\'{i}vel em: \url{https://dfrlab.org/2024/10/brazil-deepfake-challenges}.
Acesso em: 15 jun.\ 2025.

\bibitem{gao2021res2net}
GAO, S. et al.
Res2Net: A New Multi-Scale Backbone Architecture.
\textbf{IEEE Transactions on Pattern Analysis and Machine Intelligence},
v.~43, n.~2, p.~652--662, 2021.

\bibitem{gong2021ast}
GONG, Y.; CHUNG, Y.; GLASS, J.
AST: Audio Spectrogram Transformer.
In: \textbf{Proc.\ Interspeech 2021}, p.~571--575, 2021.

\bibitem{gulati2020conformer}
GULATI, A. et al.
Conformer: Convolution-augmented Transformer for Speech Recognition.
In: \textbf{Proc.\ Interspeech 2020}, p.~5036--5040, 2020.

\bibitem{heidari2024deepfake}
HEIDARI, A.; CELEBI, M.~E.
Deepfake detection using deep learning methods: A systematic and comprehensive
review.
\textbf{WIREs Data Mining and Knowledge Discovery},
v.~14, n.~2, e1520, 2024.

\bibitem{hsu2021hubert}
HSU, W.-N. et al.
HuBERT: Self-Supervised Speech Representation Learning by Masked Prediction
of Hidden Units.
\textbf{IEEE/ACM Transactions on Audio, Speech, and Language Processing},
v.~29, p.~3451--3460, 2021.

\bibitem{jung2022aasist}
JUNG, J.~W. et al.
AASIST: Audio Anti-Spoofing using Integrated Spectro-Temporal Graph Attention
Networks.
In: \textbf{IEEE International Conference on Acoustics, Speech and Signal
Processing (ICASSP)}, p.~6367--6371, 2022.

\bibitem{kameoka2018stargan}
KAMEOKA, H. et al.
StarGAN-VC: Non-parallel many-to-many voice conversion using star generative
adversarial networks.
In: \textbf{IEEE Spoken Language Technology Workshop (SLT)},
p.~266--273, 2018.

\bibitem{kurnaz2024hybrid}
KURNAZ, S.
A Hybrid CNN-LSTM Approach for Precision Deepfake Image Detection Based on
Transfer Learning.
\textbf{Electronics (MDPI)},
v.~13, n.~9, p.~1662, 2024.

\bibitem{lundberg2017unified}
LUNDBERG, S.~M.; LEE, S.-I.
A Unified Approach to Interpreting Model Predictions.
In: \textbf{Advances in Neural Information Processing Systems (NeurIPS)},
v.~30, p.~4765--4774, 2017.

\bibitem{muller2022warning}
M\"{U}LLER, N.~M. et al.
Warning: Humans cannot reliably detect speech deepfakes.
\textbf{PLoS ONE},
v.~17, n.~8, e0271552, 2022.

\bibitem{oord2016wavenet}
VAN DEN OORD, A. et al.
WaveNet: A generative model for raw audio.
\textbf{arXiv preprint} arXiv:1609.03499, 2016.

\bibitem{pindrop2025}
PINDROP.
Deepfake Trends to Look Out for in 2025.
2025.
Dispon\'{i}vel em: \url{https://www.pindrop.com/article/deepfake-trends/}.
Acesso em: 15 jun.\ 2025.

\bibitem{sahidullah2015comparison}
SAHIDULLAH, M. et al.
A Comparison of Features for Synthetic Speech Detection.
In: \textbf{Proc.\ Interspeech 2015}, p.~2087--2091, 2015.

\bibitem{shen2018natural}
SHEN, J. et al.
Natural TTS Synthesis by Conditioning WaveNet on Mel Spectrogram Predictions.
In: \textbf{IEEE International Conference on Acoustics, Speech and Signal
Processing (ICASSP)}, p.~4779--4783, 2018.

\bibitem{tak2021rawnet2}
TAK, H. et al.
End-to-End Anti-Spoofing with RawNet2.
In: \textbf{IEEE International Conference on Acoustics, Speech and Signal
Processing (ICASSP)}, p.~6369--6373, 2021.

\bibitem{todisco2016new}
TODISCO, M.; DELGADO, H.; EVANS, N.
A New Feature for Automatic Speaker Verification Anti-Spoofing:
Constant Q Cepstral Coefficients.
In: \textbf{Proc.\ Odyssey 2016 --- The Speaker and Language Recognition
Workshop}, p.~283--290, 2016.

\bibitem{todisco2019asvspoof}
TODISCO, M. et al.
ASVspoof 2019: Future Horizons in Spoofed and Fake Speech Detection.
In: \textbf{Proc.\ Interspeech 2019}, p.~1008--1012, 2019.

\bibitem{wang2023vall}
WANG, C. et al.
VALL-E: Neural Codec Language Models are Zero-Shot Text-to-Speech Synthesizers.
\textbf{arXiv preprint} arXiv:2301.02111, 2023.
Dispon\'{i}vel em: \url{https://arxiv.org/abs/2301.02111}.
Acesso em: 15 jun.\ 2025.

\bibitem{wang2024asvspoof5}
WANG, X. et al.
ASVspoof 5: Crowdsourced Speech Data, Deepfakes, and Adversarial Attacks at
Scale.
\textbf{arXiv preprint} arXiv:2408.08739, 2024.
Dispon\'{i}vel em: \url{https://arxiv.org/abs/2408.08739}.
Acesso em: 15 jun.\ 2025.

\bibitem{yamagishi2021asvspoof}
YAMAGISHI, J. et al.
ASVspoof 2021: Accelerating Progress in Spoofed and Deepfake Speech Detection.
In: \textbf{Proc.\ ASVspoof 2021 Workshop}, p.~47--54, 2021.

\bibitem{yi2022add}
YI, J. et al.
ADD 2022: The First Audio Deep Synthesis Detection Challenge.
In: \textbf{IEEE International Conference on Acoustics, Speech and Signal
Processing (ICASSP)}, p.~9216--9220, 2022.

\end{thebibliography}

% =====================================================================
%  APENDICES
% =====================================================================
\appendix

\section{Configura\c{c}\~{o}es Detalhadas do Pipeline}

\subsection{Par\^{a}metros de Pr\'{e}-processamento}

\begin{table}[H]
\centering
\caption{Configura\c{c}\~{o}es do AudioPreprocessor}
\begin{tabular}{ll}
\toprule
\textbf{Par\^{a}metro} & \textbf{Valor} \\
\midrule
Sample Rate            & 16.000 Hz \\
Dura\c{c}\~{a}o de entrada & 1,0 s (16.000 amostras) \\
Frame Length           & 25,0 ms \\
Frame Shift            & 10,0 ms \\
N\_FFT                 & 2048 \\
Hop Length             & 256 \\
Window Length          & 1024 \\
VAD                    & Energia (fallback; limiar $-25$\,dB) \\
Min.\ Dura\c{c}\~{a}o Segmento & 0,3 s \\
Pre-emphasis           & 0,97 \\
Target Loudness (AGC)  & $-23$\,LUFS \\
Faixa de Frequ\^{e}ncia & 80\,Hz -- 8.000\,Hz \\
\bottomrule
\end{tabular}
\end{table}

\subsection{Hiperpar\^{a}metros das Arquiteturas (Reais)}

\begin{table}[H]
\centering
\caption{Hiperpar\^{a}metros reais utilizados no treinamento}
\small
\begin{tabular}{lcc}
\toprule
\textbf{Par\^{a}metro} & \textbf{MultiscaleCNN} & \textbf{Ensemble Adapt.} \\
\midrule
%% AUTOGEN_BEGIN:tab_hyperparams
Learning Rate (inicial)        & 0,001 & 0,001 \\
Batch Size                     & 32 & 32 \\
L2 Regulariza\c{c}\~{a}o      & 0,0001 & 0,0001 \\
\'{E}pocas m\'{a}x.\           & 20 & 20 \\
Early Stopping (paci\^{e}ncia) & 10 & 10 \\
LR Patience                    & 5 & 5 \\
\'{E}pocas executadas          & 20 & 20 \\
Melhor val.\ acur\'{a}cia      & 96,67\% & 92,22\% \\
Par\^{a}metros                 & 48.965.581 & 2.340.650 \\
Mem\'{o}ria modelo             & 186,8 MB & 8,9 MB \\
%% AUTOGEN_END:tab_hyperparams
\bottomrule
\end{tabular}
\end{table}

\begin{table}[H]
\centering
\caption{Arquiteturas que n\~{a}o convergiram em CPU}
\small
\begin{tabular}{lccccc}
\toprule
\textbf{Arquitetura} & \textbf{Params} & \textbf{MB} &
\textbf{Lat.\,(ms)} & \textbf{\'{E}poc.} & \textbf{Causa} \\
\midrule
%% AUTOGEN_BEGIN:tab_nao_convergiu
EfficientNet-LSTM & 7.986.852 & 30,5 & 63,8 & 11 & Seq.\ PCM longa \\
AASIST & 1.150.592 & 4,4 & 89,0 & 11 & SincConv+GAT, CPU \\
RawNet2 & 5.303.553 & 20,2 & 156,4 & 11 & GRU sobre PCM, CPU \\
%% AUTOGEN_END:tab_nao_convergiu
\bottomrule
\end{tabular}
\end{table}

\section{Implementa\c{c}\~{a}o de Refer\^{e}ncia (TensorFlow/Keras)}

\subsection{Pipeline Principal}

\begin{lstlisting}[language=Python,
caption={Estrutura Principal do Pipeline (TensorFlow/Keras)}]
import numpy as np
import librosa
import tensorflow as tf
from app.models import ModelFactory
from app.features import FeatureExtractor

class DeepAudioPipeline:
    def __init__(self, architecture='multiscale_cnn'):
        self.feat_ext = FeatureExtractor(sample_rate=16000)
        self.model = ModelFactory.create(architecture)

    def predict(self, audio_path: str) -> dict:
        # Pre-processamento
        wav, _ = librosa.load(audio_path, sr=16000, mono=True)
        wav = self._apply_agc(wav, target_lufs=-23)
        wav = self._apply_vad(wav, sr=16000, threshold_db=-25)

        # Extracao de features (espectrograma mel)
        mel = librosa.feature.melspectrogram(
            y=wav, sr=16000, n_mels=80,
            n_fft=2048, hop_length=256
        )
        mel_db = librosa.power_to_db(mel, ref=np.max)

        # Classificacao
        x = tf.expand_dims(
            tf.expand_dims(mel_db.T, axis=0), axis=-1
        )
        score = float(self.model.predict(x)[0][0])

        return {
            "prediction": "FAKE" if score > 0.5 else "REAL",
            "confidence": score,
            "model": self.model.name
        }
\end{lstlisting}

\subsection{Treinamento com Early Stopping}

\begin{lstlisting}[language=Python,
caption={Loop de Treinamento com Early Stopping e ReduceLROnPlateau}]
import tensorflow as tf

callbacks = [
    tf.keras.callbacks.EarlyStopping(
        monitor='val_accuracy',
        patience=10,
        restore_best_weights=True,
        verbose=1
    ),
    tf.keras.callbacks.ReduceLROnPlateau(
        monitor='val_loss',
        factor=0.5,
        patience=5,
        min_lr=1e-6,
        verbose=1
    ),
    tf.keras.callbacks.ModelCheckpoint(
        filepath='results/models/{arch}_best.h5',
        monitor='val_accuracy',
        save_best_only=True
    )
]

model.compile(
    optimizer=tf.keras.optimizers.Adam(learning_rate=1e-3),
    loss='binary_crossentropy',
    metrics=['accuracy']
)

history = model.fit(
    X_train, y_train,
    validation_data=(X_val, y_val),
    epochs=20,
    batch_size=32,
    callbacks=callbacks
)
\end{lstlisting}

\section{Resultados Completos dos Experimentos}

\subsection{Resumo Consolidado}

\begin{table}[H]
\centering
\caption{Resultados Experimentais Consolidados}
\small
\begin{tabular}{lcccccc}
\toprule
\textbf{Arquitetura} & \textbf{Acur.} & \textbf{EER} & \textbf{AUC} &
\textbf{Lat.} & \textbf{MB} & \textbf{Status} \\
\midrule
%% AUTOGEN_BEGIN:tab_consolidado
MultiscaleCNN & 97,78\% & 5,56\% & 0,997 & 133,7ms & 186,8 & OK \\
Ensemble Adapt. & 91,11\% & 13,33\% & 0,975 & 48,1ms & 8,9 & OK \\
Effic.-LSTM & 50,00\% & 46,67\% & 0,511 & 63,8ms & 30,5 & n/a \\
AASIST & 50,00\% & 50,00\% & 0,500 & 89,0ms & 4,4 & n/a \\
RawNet2 & 50,00\% & 50,00\% & 0,000 & 156,4ms & 20,2 & n/a \\
%% AUTOGEN_END:tab_consolidado
\bottomrule
\end{tabular}
\end{table}

\subsection{Robustez (Acur\'{a}cia \% por SNR)}

\begin{table}[H]
\centering
\caption{Dados completos de robustez --- AUC-ROC complementar}
\begin{tabular}{lcccc}
\toprule
\textbf{Arquitetura} & \textbf{Limpo} & \textbf{SNR 30dB} &
\textbf{SNR 20dB} & \textbf{SNR 10dB} \\
\midrule
%% AUTOGEN_BEGIN:tab_robustez_completo
\multicolumn{5}{l}{\textit{Acur\'{a}cia (\%)}} \\
MultiscaleCNN & 97,78 & 71,11 & 72,22 & 58,89 \\
Ensemble Adapt. & 87,78 & 74,44 & 50,00 & 50,00 \\
\midrule
\multicolumn{5}{l}{\textit{EER (\%)}} \\
MultiscaleCNN & 2,22 & 28,89 & 27,78 & 41,11 \\
Ensemble Adapt. & 12,22 & 25,56 & 50,00 & 50,00 \\
\midrule
\multicolumn{5}{l}{\textit{AUC-ROC}} \\
MultiscaleCNN & 0,997 & 0,994 & 0,808 & 0,652 \\
Ensemble Adapt. & 0,966 & 0,923 & 0,748 & 0,714 \\
%% AUTOGEN_END:tab_robustez_completo
\bottomrule
\end{tabular}
\end{table}

\end{document}
